% Arbeitsblatt Mathematik Klasse 10 (German, Gymnasium), with a tikz sketch and a pgfplots parabola.
% Headings: "Aufgabe N (P BE)" in bold (BE = Bewertungseinheiten, the usual school unit); parts "a)", "b)", ...
\documentclass[12pt,a4paper]{article}
\usepackage[T1]{fontenc}
\usepackage[utf8]{inputenc}
\usepackage{lmodern}
\usepackage[ngerman]{babel}
\usepackage[margin=2.3cm,top=1.8cm,bottom=1.9cm]{geometry}
\usepackage{amsmath,amssymb}
\usepackage[locale=DE]{siunitx}
\DeclareSIUnit{\euro}{\text{€}}
\usepackage{tikz}
\usetikzlibrary{patterns,arrows.meta}
\usepackage{pgfplots}
\pgfplotsset{compat=1.18}
\usepackage{enumitem}

\setlength{\parindent}{0pt}
\setlist[enumerate,1]{label=\alph*), leftmargin=2em, itemsep=0.3em, topsep=0.3em}
\newcommand{\aufgabe}[2]{\bigskip\par\textbf{Aufgabe #1 (#2 BE)}\par\nopagebreak\smallskip}

\begin{document}

% ---- Kopf -----------------------------------------------------------------------------
{\small Städtisches Gymnasium \textperiodcentered{} Mathematik \textperiodcentered{} Klasse 10 \hfill Schuljahr 2026/27}\\[-0.4em]
\rule{\linewidth}{0.6pt}\\[0.5em]
{\Large\bfseries Arbeitsblatt 4: Quadratische Gleichungen, Pythagoras, Prozent- und Zinsrechnung}\\[0.5em]
Name: \rule{5cm}{0.4pt} \hfill Datum: \rule{3cm}{0.4pt} \hfill Gesamt: 30 BE\\[0.3em]
{\small Taschenrechner ist erlaubt. Schreibe deinen Lösungsweg vollständig auf.}

% ---- Aufgabe 1 ------------------------------------------------------------------------
\aufgabe{1}{6}
Löse die quadratischen Gleichungen mit der pq-Formel. Gib die Lösungsmenge an.
\begin{enumerate}
  \item $x^2 - 6x + 8 = 0$
  \item $x^2 + 4x - 21 = 0$
  \item $2x^2 - 12x + 18 = 0$
  \item $x^2 + 2x + 5 = 0$
\end{enumerate}

% ---- Aufgabe 2 ------------------------------------------------------------------------
\aufgabe{2}{7}
Eine Leiter der Länge \qty{5.0}{m} lehnt an einer senkrechten Hauswand. Ihr Fuß steht \qty{1.4}{m}
von der Wand entfernt (siehe Skizze).

\begin{center}
\begin{tikzpicture}[scale=0.95]
  \draw[thick] (-0.3,0) -- (4.6,0);
  \fill[pattern=north east lines, pattern color=gray] (-0.3,-0.2) rectangle (4.6,0);
  \draw[very thick] (0,0) -- (0,4.7);
  \draw[very thick, brown!70!black] (1.26,0) -- (0,4.32);
  \draw (0,0.25) -- (0.25,0.25) -- (0.25,0);
  \draw[<->] (0,-0.45) -- (1.26,-0.45) node[midway, below] {\qty{1,4}{m}};
  \node[above right] at (0.75,2.45) {$l = \qty{5,0}{m}$};
  \node[left] at (0,2.2) {$h$};
  \draw (0.85,0) arc[start angle=180, end angle=105.5, radius=0.4];
  \node at (0.85,0.55) {$\alpha$};
\end{tikzpicture}

\small Skizze (nicht maßstabsgetreu)
\end{center}
\begin{enumerate}
  \item Wie hoch reicht die Leiter an der Wand?
  \item Der Fuß der Leiter soll so weit von der Wand entfernt stehen, dass die Leiter noch \qty{4.0}{m}
        hoch reicht. Wie groß ist dieser Abstand?
  \item Berechne den Winkel $\alpha$, den die Leiter mit dem Boden bildet (Skizze oben).
\end{enumerate}

% ---- Aufgabe 3 ------------------------------------------------------------------------
\aufgabe{3}{8}
Prozent- und Zinsrechnung.
\begin{enumerate}
  \item Ein Fahrrad kostet \qty{480}{\euro}. Im Schlussverkauf wird der Preis um $15\,\%$ gesenkt. Wie viel
        kostet das Fahrrad jetzt?
  \item Nach dem Schlussverkauf wird der reduzierte Preis wieder um $15\,\%$ erhöht. Kostet das Fahrrad nun
        wieder \qty{480}{\euro}? Begründe deine Antwort mit einer Rechnung.
  \item Frau Weber legt \qty{2500}{\euro} für 4 Jahre mit einem Zinssatz von $3\,\%$ pro Jahr an. Die Zinsen
        werden jedes Jahr mitverzinst. Wie hoch ist ihr Guthaben nach 4 Jahren?
  \item Wie viel Zinsen bringen \qty{800}{\euro} bei $2{,}5\,\%$ Jahreszinsen in 90 Tagen? Rechne mit
        360 Tagen pro Jahr.
\end{enumerate}

% ---- Aufgabe 4 ------------------------------------------------------------------------
\aufgabe{4}{9}
Die Abbildung zeigt den Graphen einer quadratischen Funktion $f$.

\begin{center}
\begin{tikzpicture}
\begin{axis}[width=8.4cm, height=7cm, axis lines=middle,
    xlabel={$x$}, ylabel={$y$},
    xmin=-3, xmax=5, ymin=-5, ymax=6,
    xtick={-3,-2,...,5}, ytick={-5,-4,...,6}, grid=both, grid style={gray!25},
    tick label style={font=\footnotesize},
    every axis x label/.style={at={(ticklabel* cs:1.0)}, anchor=west},
    every axis y label/.style={at={(ticklabel* cs:1.0)}, anchor=south}]
  \addplot[thick, domain=-1.8:3.8, samples=70] {x^2-2*x-3};
\end{axis}
\end{tikzpicture}
\end{center}
\begin{enumerate}
  \item Lies den Scheitelpunkt und die Nullstellen von $f$ aus dem Graphen ab.
  \item Gib die Funktionsgleichung in Scheitelpunktform an und forme sie in die Normalform
        $f(x) = x^2 + px + q$ um.
  \item Für welche $x$-Werte gilt $f(x) < 0$?
  \item Der Punkt $P(4\mid y)$ liegt auf dem Graphen von $f$. Berechne $y$.
\end{enumerate}

\end{document}
